\begin{frame}{Reference: input and source editing}
\small
\begin{tabularx}{\textwidth}{@{}P{54mm}Y@{}}
\toprule
\textbf{Input} & \textbf{Behavior}\\
\midrule
\cmd{10 source} & Insert/replace; conservative final-semicolon convenience\\
\cmd{10: exact source} & Literal physical line; no semicolon insertion or continuation\\
\cmd{10} & Delete that line\\
\cmd{DEL range}, \cmd{DELETE range} & Delete an inclusive range\\
\cmd{RENUM [start [step]]} & Renumber; defaults 10 and 10; REORDER is an alias\\
\cmd{UNDO} & Restore one of up to 50 source changes\\
\cmd{NEW}, \cmd{CLEAR} & Clear source and filename after unsaved-source guard\\
\cmd{STATUS} & Source count, dirty status, filename, active AUTO settings\\
\bottomrule
\end{tabularx}
\captionline{Editing numbers: 1 through 999999999. EDIT, IMPORT, and LOAD preserve literal source text.}
\note{Reference only. Semicolon insertion is a typing convenience, not a source rewriter or a full parser. For exact punctuation use the colon form. Empty literal numbered lines can preserve a blank physical line. A normal number with no code deletes. UNDO restores source only and does not restore an earlier filename or undo writes.}
\end{frame}

\begin{frame}{Reference: AUTO and explicit evaluation}
\small
\begin{tabularx}{\textwidth}{@{}P{51mm}Y@{}}
\toprule
\textbf{Input} & \textbf{Behavior}\\
\midrule
\cmd{AUTO [start [step]]} & Default start highest + 5, or 5; default step 5; skips occupied numbers\\
\cmd{AUTO ECHO ON|OFF} & Control numbered duplicates; independent of AUTO activation\\
\cmd{AUTO OFF} & End automatic storage\\
\cmd{: exact source} & Literal one-line entry while AUTO is active\\
\cmd{EVAL code} & Immediate execution, uppercase only; leaves AUTO active\\
\cmd{? expression} & Scalar evaluation and terse Dumper display\\
\key{Ctrl-C}, \key{Ctrl-D} in AUTO & Cancel unfinished input and end AUTO; keep completed source\\
\bottomrule
\end{tabularx}
\captionline{Multiline AUTO groups preserve their own semicolons and store as one undoable source edit.}
\note{Reference only. In plain terminal mode Ctrl-D should be typed at an empty physical input line. Blank input does nothing. Commands remain recognized during AUTO. Lowercase eval is ordinary Perl, so AUTO stores it.}
\end{frame}

\begin{frame}[fragile]{Reference: EVAL and Perl eval}
\begin{lstlisting}[style=perl]
EVAL say "immediate, even during AUTO"

eval { die "example\n" }; say "Caught: $@";
eval 'say "string eval"';
\end{lstlisting}
\par\vspace{2mm}
Uppercase \cmd{EVAL} is an AnteREPL command.\par\medskip
Lowercase \cmd{eval} retains Perl's exception-catching and string-eval behavior.\par\medskip
Read \cmd{\$@} in the same entry; a later REPL evaluation can change it.
\note{Reference only. Other prompt commands are case insensitive. An uppercase EVAL prefix removes ambiguity when a line begins with an editing number or resembles a command. EVAL does not automatically print an expression's value; use ? for that. These examples are alternatives, not instructions to paste as a single entry while AUTO is active.}
\end{frame}

\begin{frame}{Reference: LIST and display controls}
\small
\begin{tabularx}{\textwidth}{@{}P{51mm}Y@{}}
\toprule
\textbf{Command or key} & \textbf{Behavior}\\
\midrule
\cmd{LIST}, \cmd{LIST 10} & All source, or one editing number\\
\cmd{LIST 10-50}, \cmd{LIST10-50} & Inclusive range, with optional compact spelling\\
\cmd{LIST 10-}, \cmd{LIST -50} & Open-ended inclusive ranges\\
\cmd{LIST [range] SLOW} & Start controlled listing slowly\\
\cmd{LIST [range] PAUSE} & Pause after each screenful\\
Space / S / Enter & Pause or resume / toggle speed / step one row\\
Q / Esc / Ctrl-C & Stop the listing\\
Other typing or movement & Stop LIST and pass the key to the editor\\
\bottomrule
\end{tabularx}
\captionline{Normal controlled speed is about 33 lines/second; slow speed about 4. Plain or piped output never waits for keys.}
\note{Reference only. Numbers are left aligned within the widest selected number's width, with right padding so source begins in the same column. Short normal listings finish immediately. LIST modifies only the display until a displayed numbered row is edited and submitted.}
\end{frame}

\begin{frame}{Reference: file commands}
\small
\begin{tabularx}{\textwidth}{@{}P{49mm}Y@{}}
\toprule
\textbf{Command} & \textbf{Behavior}\\
\midrule
\cmd{LOAD file} & Read plain or native source; remember destination and format\\
\cmd{IMPORT file} & Read plain source as an unsaved copy; no adopted filename\\
\cmd{SAVE [file]} & Save source; .pl selects plain, .bperl selects native\\
\cmd{EXPORT file} & Write source plus default pragmas; does not mark source saved\\
\cmd{EDIT} & Editor round trip; changed source renumbers 10, 20, \ldots\\
\cmd{LOAD!}, \cmd{IMPORT!} & Deliberately discard unsaved source before reading\\
\cmd{NEW!}, \cmd{CLEAR!} & Deliberately discard unsaved source and filename\\
\cmd{QUIT!}, \cmd{EXIT!} & Exit despite unsaved source\\
\bottomrule
\end{tabularx}
\captionline{UTF-8 source. CRLF input becomes LF. Invalid native records, duplicates, NUL, and bare CR are rejected.}
\note{Reference only. Native .bperl requires the PERL ANTE REPL header; legacy BASIC PERL V0.1 headers remain readable. Either header also identifies native input under another extension. SAVE and EXPORT replace the named target using a temporary file and rename. A GUI editor must wait for editing to finish. Failure or invalid output leaves source unchanged. An unchanged EDIT preserves numbers. LOAD and EDIT do not syntax-check or execute source.}
\end{frame}

\begin{frame}{Reference: environment operations}
\small
\begin{tabularx}{\textwidth}{@{}P{56mm}Y@{}}
\toprule
\textbf{Command} & \textbf{Behavior}\\
\midrule
\cmd{ENV [name]} & Sorted VAR=value output for all entries or one name\\
\cmd{ENV SET name value} & Set a literal value; quote spaces or an empty string\\
\cmd{ENV UNSET name} & Remove a variable\\
\cmd{EDIT ENV}, \cmd{ENV EDIT} & Edit the whole live environment; deleted lines unset names\\
\cmd{ENV SAVE file} & Write a complete snapshot\\
\cmd{ENV LOAD file} & Replace the whole environment from a snapshot\\
\cmd{ENV RELOAD [file]} & Read a named or last saved/loaded environment file\\
\cmd{ENV RESET} & Restore the environment inherited at startup\\
\bottomrule
\end{tabularx}
\captionline{No environment UNDO. Failed parsing or editing leaves the current environment intact.}
\note{Reference only. These operations affect the running process and future child commands, not the parent shell. Direct assignments and deletes in percent-ENV work too. Subcommand names such as SET are reserved after ENV; inspect a variable with that name using dump \$ENV{SET}. A snapshot may contain secrets; do not commit real snapshots or include them in demo material.}
\end{frame}

\begin{frame}[fragile]{Reference: VAR=value file syntax}
\begin{lstlisting}
# A small example, not a complete process snapshot
EDITOR=vi
GREETING="hello world"
LITERAL='$HOME stays literal'
MULTILINE="line one\nline two"
export EXAMPLE_COLOR=blue
\end{lstlisting}
\par\vspace{2mm}
Blank lines and comments are allowed. Single quotes preserve literal text.\par\medskip
Double quotes support escapes. No variable or command expansion occurs.\par\medskip
Duplicate names and NUL bytes make the entire replacement fail.
\note{Reference only. Do not LOAD this small sample casually: environment LOAD replaces, rather than merges, the process environment. A whitespace-preceded hash begins an inline comment, while abc\#def preserves the hash. Supported double-quoted escapes are newline, carriage return, tab, backslash, double quote, and hex byte. Names cannot contain whitespace, equals, hash, quotes, or NUL.}
\end{frame}

\begin{frame}{Reference: screen keys}
\small
\begin{tabularx}{\textwidth}{@{}P{48mm}Y@{}}
\toprule
\textbf{Key} & \textbf{Behavior}\\
\midrule
Arrows; Left at row start & Traverse screen; Left crosses to previous row's end\\
Enter & Submit selected row or its complete multiline group\\
Ctrl-N / Shift+Enter & Insert newline without submitting; Shift+Enter needs terminal support\\
Insert; Backspace / Delete & Toggle insert/overwrite; remove characters\\
Home / End; PageUp / PageDown & Row ends; traverse retained screen rows\\
Esc / Ctrl-End; Ctrl-Home & Fresh bottom row; oldest retained row\\
Ctrl-U / Ctrl-L & Clear row / redraw screen\\
Tab & Spaces to next four-column tab stop\\
Ctrl-C / Ctrl-D & Cancel or interrupt / end AUTO, otherwise guarded exit\\
\bottomrule
\end{tabularx}
\note{Reference only. Default editing mode is insert. Left stops at the oldest retained row and visits empty rows one by one. Right currently stops at the end of its row. Shift+Enter is recognized as CSI-u ESC [ 13 ; 2 u or xterm ESC [ 27 ; 2 ; 13 tilde, without the spaces. No app can distinguish Shift+Enter when the terminal sends ordinary Enter for both. Ctrl-N is the fallback. Input to a running program has its own editing boundary.}
\end{frame}

\begin{frame}[fragile]{Reference: functions and old code references}
\begin{lstlisting}
EVAL sub triple { return 3 * shift; }
$old = \&triple
EVAL sub triple { return 4 * shift; }
(*@\result{Subroutine triple redefined ...}@*)
say $old->(14)
(*@\result{42}@*)
say triple(14)
(*@\result{56}@*)
\end{lstlisting}
\captionline{Assumes a fresh workspace. FUNCTIONS includes live imported names and aliases, but not anonymous subs held only in variables or bodyless stubs.}
\note{Reference only. A code reference captured before redefinition keeps the earlier body. no warnings 'redefine' in the submitted entry can suppress the normal warning. FUNCTIONS lists only the current package, not every package in the process. RESET selects another workspace without recursively erasing the old package graph, so do not use it as a memory-reclamation or process-isolation guarantee.}
\end{frame}

\begin{frame}{Reference: startup, help, shell, and exit}
\small
\begin{tabularx}{\textwidth}{@{}P{59mm}Y@{}}
\toprule
\textbf{Invocation or command} & \textbf{Behavior}\\
\midrule
\cmd{anterepl [program.pl]} & Start fresh or load a file without running it\\
\cmd{--quiet}, \cmd{--plain} & Hide banner; disable screen mode\\
\cmd{--no-readline} & Alias for --plain\\
\cmd{--help}, \cmd{-h}, \cmd{--version} & Print HELP or application version and exit\\
\cmd{-- -example.pl} & Accept startup filename beginning with a hyphen\\
\cmd{! command}, \cmd{!! command} & Captured shell output; interactive terminal handoff\\
\cmd{HELP} & Complete built-in reference\\
\cmd{QUIT}, \cmd{EXIT} & Exit unless source is unsaved; add ! to discard it\\
\bottomrule
\end{tabularx}
\captionline{Startup also accepts .bperl files. Outside AUTO, Ctrl-D is a guarded exit.}
\note{Reference only. Piped input uses plain mode and errors or unfinished/unsaved input produce a nonzero exit. Native Windows shell syntax differs from Unix. In plain mode both shell escapes use normal terminal I/O. Perl's own exit terminates the app directly; prefer QUIT when you want the source guard.}
\end{frame}

\begin{frame}{Reference: loading an existing Perl script}
\large
\cmd{LOAD} imports source into the current interpreter.\par\vspace{4mm}
\normalsize
Check scripts that rely on process startup, \cmd{\_\_DATA\_\_}, unusual compile-time context, or external terminal control.\par\medskip
\cmd{RUN} shares workspace state and applies the evaluator's default pragmas.\par\medskip
For ordinary script process semantics, use \cmd{!! perl filename.pl}.\par\medskip
For stored inspection code, call \cmd{Data::Dumper::Dumper} or \cmd{Devel::Peek::Dump} explicitly.
\note{Reference only. The source mechanism is intentionally small, not a general loader that emulates every property of perl script.pl. A string eval does not create a normal script DATA handle or a new interpreter lifecycle. Source can use packages and modules, but code that names another package is outside the default workspace inventory. The HELP and POD provide the authoritative current limitations.}
\end{frame}

\begin{frame}{Sources and version of this talk}
\small
Implementation snapshot: \cmd{6a9c90f7f189} \textemdash{} November 18, 2026.\par\medskip
\href{https://github.com/oodler577/p5-App-AnteREPL/blob/6a9c90f7f189907b6c75911871b015fd8cb5582d/lib/App/AnteREPL.pm}{\textbf{Implementation, complete HELP, and POD}}\par\medskip
\href{https://github.com/oodler577/p5-App-AnteREPL/wiki}{\textbf{User guide and workflow documentation}}\par\medskip
\href{https://github.com/oodler577/p5-App-AnteREPL/tree/6a9c90f7f189907b6c75911871b015fd8cb5582d/t}{\textbf{Unit and screen tests}}\hspace{5mm}
\href{https://github.com/oodler577/p5-App-AnteREPL/blob/6a9c90f7f189907b6c75911871b015fd8cb5582d/tools/terminal-smoke.pl}{\textbf{Perl terminal harness}}\par\medskip
\href{https://github.com/oodler577/p5-App-AnteREPL/blob/6a9c90f7f189907b6c75911871b015fd8cb5582d/CONTRIBUTING.md}{\textbf{Distribution and release notes}}
\par\vspace{4mm}
This deck documents an implementation under active UX testing.\par
The included README explains rehearsal, shorter routes, and Overleaf editing.
\note{Sources are the project's implementation and its user-facing documentation, not reconstructed from historical chat alone. Wiki snapshot 8b190d3 documents the same Left-navigation behavior. Examples use generic data and the application's existing triple demonstration. Update the implementation hash and any changed behavior together when revising the talk. No external artwork, logos, or C64 screen photographs are needed for this text-based deck.}
\end{frame}
